% =============================================================================
% Veridian: IEEE BDePIN-Energy 2026 Workshop Paper
% IEEE Two-Column Conference Format (IEEEtran)
% =============================================================================
\documentclass[conference]{IEEEtran}
\IEEEoverridecommandlockouts

% --- Packages ---
\usepackage{cite}
\usepackage{amsmath,amssymb,amsfonts}
\usepackage{graphicx}
\usepackage{textcomp}
\usepackage{xcolor}
\usepackage{booktabs}
\usepackage{array}
\usepackage{url}
\usepackage{hyperref}
\hypersetup{
  colorlinks=true,
  linkcolor=black,
  citecolor=black,
  urlcolor=black
}

\def\BibTeX{{\rm B\kern-.05em{\sc i\kern-.025em b}\kern-.08em
    T\kern-.1667em\lower.7ex\hbox{E}\kern-.125emX}}

% =============================================================================
\begin{document}
% =============================================================================

\title{Veridian: A Generalized Cryptoeconomic Lifecycle Architecture for Decentralized Energy Assets}

\author{%
\IEEEauthorblockN{Kannan Sasinthiran}
\IEEEauthorblockA{\textit{Dept.\ of Statistics and Data Science}\\
National University of Singapore\\
Singapore\\
e1121829@u.nus.edu}
\and
\IEEEauthorblockN{Alican Ala\c{s}\i{}k}
\IEEEauthorblockA{\textit{UZH Blockchain Center}\\University of Zurich\\
Zurich, Switzerland\\
alican.alasik@uzh.ch}
\and
\IEEEauthorblockN{Tao}
\IEEEauthorblockA{\textit{UZH Blockchain Center}\\University of Zurich\\
Zurich, Switzerland}
}

\maketitle

% =============================================================================
\begin{abstract}
Decentralized Physical Infrastructure Networks (DePIN) promise to democratize the deployment of hardware infrastructure by aligning the economic interests of investors, operators, and consumers. However, existing energy DePIN protocols struggle to scale into regulated environments due to an absence of formalized trust boundaries between physical hardware, off-chain legal compliance, and on-chain settlement. 

This paper presents a generalized Cryptoeconomic Lifecycle Architecture for distributed energy assets. We formalize this lifecycle into five distinct stages: Capital Formation, Compliance-Gated Ownership, Attested Telemetry, Deterministic Settlement, and Governance Constraints. For each stage, we provide an explicit mapping of where security is guaranteed by cryptography, economics, law, or plain trust.

To demonstrate the viability of this architecture, we instantiate it via \emph{Veridian}, an Equipment-as-a-Service (EaaS) prototype for retail solar infrastructure in South Africa. The implementation utilizes a 3-of-5 multisignature escrow for capital formation, ERC-3643 tokens for compliance, TPM 2.0 hardware for physical telemetry, and a deterministic 75/8/7/5/5 fee waterfall executing on Optimism Sepolia. An end-to-end simulation of a 450~W solar panel yields a verifiable net investor return of 18.7\% per annum. Finally, we establish the generalizability of the architecture by extending it to peer-to-peer excess battery storage pooling for localized data centers.
\end{abstract}

\begin{IEEEkeywords}
blockchain, decentralized physical infrastructure, DePIN, energy assets, equipment-as-a-service, ERC-3643, TPM 2.0, real-world assets
\end{IEEEkeywords}

% =============================================================================
\section{Introduction}
% =============================================================================

The integration of blockchain technology into the energy sector has primarily focused on peer-to-peer (P2P) trading of generated capacity \cite{pop2018peer,andoni2019blockchain}. More recently, Decentralized Physical Infrastructure Networks (DePIN) have emerged as an architectural pattern to coordinate the deployment of the physical hardware itself through on-chain incentive structures \cite{lin2024depin}. 

Applied to energy infrastructure, DePIN protocols can theoretically align the economic interests of equipment investors, field operators, and energy consumers without a central coordinating entity. However, existing energy DePIN protocols (e.g., Power Ledger, SolarCoin) exhibit a critical weakness: they often conflate the trust assumptions required at different stages of the asset lifecycle. They rely on software-reported metrics or grid-metered totals (an unresolved oracle problem at the physical layer) while utilizing permissionless tokens that preclude regulated, compliance-gated fractional investment.

To mature from localized prototypes into globally scalable infrastructure financing systems, DePIN protocols require an honest, explicit map of trust boundaries. We argue that the contribution required by the field is not simply another tokenization platform, but rather a unified framework that accounts for which stage of an asset's life is secured by cryptography, which by economics, which by law, and which by plain trust.

This paper addresses this gap by proposing a generalized \emph{Cryptoeconomic Lifecycle Architecture} for decentralized energy assets. The primary contributions are:

\begin{enumerate}
  \item \textbf{The Lifecycle Architecture and Trust Boundary Map.} We define a five-stage framework (Capital Formation, Compliance-Gated Ownership, Attested Telemetry, Deterministic Settlement, and Governance Constraints) and explicitly map the trust dependencies across physical, legal, and on-chain execution layers.
  \item \textbf{A Reference Instantiation (Veridian).} We validate the architecture through the implementation and economic simulation of Veridian, a prototype applying this lifecycle to solar Equipment-as-a-Service (EaaS) financing in South Africa.
\end{enumerate}

The remainder of this paper is structured as follows. Section~\ref{sec:architecture} introduces the generalized Cryptoeconomic Lifecycle Architecture. Section~\ref{sec:instantiation} details the Veridian solar instantiation. Section~\ref{sec:economic} presents the economic model and simulation results. Section~\ref{sec:generalizability} explores the generalizability of the architecture to advanced use-cases, specifically battery storage pooling for data centers. Section~\ref{sec:evaluation} evaluates the trust model and compares it to existing protocols. Section~\ref{sec:conclusion} concludes.

% =============================================================================
\section{The Cryptoeconomic Lifecycle Architecture}
\label{sec:architecture}
% =============================================================================

The deployment, operation, and financing of any decentralized physical infrastructure asset follows a consistent lifecycle. By isolating this lifecycle into five discrete stages, we can explicitly define the mechanism securing each stage. Fig.~\ref{fig:architecture} illustrates this lifecycle and maps the corresponding trust boundaries.

\begin{figure}[t]
\centering
\fbox{\parbox{0.9\columnwidth}{\centering\small
\textbf{1. Capital Formation}\\
Mechanism: Threshold Escrow $\rightarrow$ \textit{Cryptography / Law}\\[4pt]
\textbf{2. Compliance-Gated Ownership}\\
Mechanism: Identity Registry (ERC-3643) $\rightarrow$ \textit{Law / Trust}\\[4pt]
\textbf{3. Attested Telemetry}\\
Mechanism: TPM 2.0 Hardware Root $\rightarrow$ \textit{Cryptography bounded by Physics}\\[4pt]
\textbf{4. Deterministic Settlement}\\
Mechanism: Smart Contract Waterfall $\rightarrow$ \textit{Cryptography}\\[4pt]
\textbf{5. Governance Constraints}\\
Mechanism: Oracle-referenced Price Caps $\rightarrow$ \textit{Economics / Trust}
}}
\caption{The five stages of the Cryptoeconomic Lifecycle Architecture and their corresponding trust boundary mappings.}
\label{fig:architecture}
\end{figure}

\subsection{Stage 1: Capital Formation}
Before physical infrastructure can be deployed, retail or institutional capital must be aggregated securely. In a decentralized environment, this requires a threshold escrow mechanism (e.g., an $M$-of-$N$ multisignature smart contract) to prevent unilateral capital flight by the platform operator. Trust at this stage is hybrid: \emph{cryptography} secures the funds on-chain, while \emph{law} governs the off-chain transfer of the physical title to the Special Purpose Vehicle (SPV) once the threshold is met.

\subsection{Stage 2: Compliance-Gated Ownership}
Regulated real-world assets (RWAs) cannot be represented by permissionless bearer instruments (like ERC-20 tokens). The ownership layer must strictly enforce Know Your Customer (KYC) and Anti-Money Laundering (AML) restrictions at the time of transfer. Trust here is anchored in \emph{law and plain trust} (relying on a regulated off-chain KYC provider), which is then strictly enforced on-chain via identity registry smart contracts.

\subsection{Stage 3: Attested Telemetry}
The defining challenge of DePIN is verifying that physical work was actually performed. The architecture dictates that telemetry must be cryptographically signed at the hardware level (e.g., via a Trusted Platform Module, TPM 2.0) before transmission. Trust is shifted from the network operator to \emph{cryptography bounded by physical hardware}—tampering with the data requires physical extraction of the secured private key \cite{boudguiga2017iot,tpm20spec}.

\subsection{Stage 4: Deterministic Settlement}
Once telemetry is verified, the resulting revenue must be distributed to stakeholders (investors, operators, platform) without custodial intermediation. A predefined execution layer contract ingests the telemetry and deterministically allocates yields via a fee waterfall. This stage is secured entirely by \emph{cryptography} (the immutable rules of the EVM or equivalent execution environment).

\subsection{Stage 5: Governance Constraints}
Physical infrastructure directly impacts human end-users. Unconstrained decentralized autonomous organizations (DAOs) optimizing purely for yield could impose extractive pricing on consumers. Therefore, the architecture requires automated on-chain circuit breakers that limit governance actions based on external oracles. This stage relies on \emph{economics} (incentive alignment) and \emph{plain trust} in the oracle data provider.

% =============================================================================
\section{Instantiation: Veridian Solar EaaS}
\label{sec:instantiation}
% =============================================================================

To validate the architecture, we implemented \emph{Veridian}, a prototype targeting the deployment of solar photovoltaics (PV) in South Africa. The country faces severe electricity access challenges, with a history of rolling load-shedding and 14.7~million citizens lacking reliable access \cite{nersa2023,worldbank2023}. While the solar resource is abundant (averaging 5.5 sun-hours per day) \cite{saws2023solar}, the upfront capital cost of a 450~W system (R6,750) is prohibitive for median-income households. The Veridian protocol addresses this through an Equipment-as-a-Service (EaaS) model, mapping directly to the five lifecycle stages.

\subsection{Capital Formation: 3-of-5 Multisignature Escrow}
In Veridian, investor USDC is aggregated in an on-chain \texttt{MultiSigEscrow} contract on Optimism Sepolia. Capital release requires 3-of-5 signatures from the Platform Operator, Field Operator, SPV Representative, Legal Representative, and Independent Verifier. Capital is only released once the Legal Representative confirms off-chain title transfer, directly linking the cryptographic state to legal ownership.

\subsection{Compliance: ERC-3643 Tokenization}
To satisfy South African Financial Sector Conduct Authority (FSCA) constraints, ownership is tokenized using the ERC-3643 standard \cite{erc3643}. The \texttt{RWAToken} contract intercepts all transfers via a \texttt{canTransfer} hook, querying an \texttt{IdentityRegistry} to ensure the recipient wallet possesses a valid KYC claim issued by a regulated entity.

\subsection{Telemetry: Hardware Root of Trust}
Each solar installation is provisioned with a simulated TPM~2.0 metering unit operating on the SECP256R1 curve. The meter produces a JSON payload containing cumulative energy (kWh) and a timestamp, signed via the TPM's private key. The off-chain yield engine verifies this signature before computing revenue. 

\subsection{Settlement: The 75/8/7/5/5 Waterfall}
\label{sec:waterfall}
Revenue ($G$) is computed as energy consumed ($q$) multiplied by the EaaS rate ($p$). The \texttt{YieldDistributor} contract deterministically splits this revenue into five tranches upon receiving an authorized oracle report:
\begin{itemize}
    \item \textbf{75\%: Token Holders} (Investor yield, with a 15\% penalty applied to unstaked tokens to encourage long-term capital commitment).
    \item \textbf{8\%: Operations and Maintenance} (Field operator).
    \item \textbf{7\%: Insurance Reserve} (Held by SPV for asset replacement).
    \item \textbf{5\%: Pool Expansion} (For acquiring new arrays).
    \item \textbf{5\%: Platform Operations} (Veridian protocol).
\end{itemize}

\subsection{Governance: Oracle-Referenced Price Cap}
To prevent the token-holder DAO from imposing monopolistic EaaS rates on households, a \texttt{PricingController} contract tracks the Eskom residential tariff ($P_{\text{grid}}$) \cite{eskom2024tariff} via an off-chain oracle. It strictly enforces $P_{\text{proposed}} \leq 0.90 \times P_{\text{grid}}$. Any DAO vote exceeding this 10\% discount ceiling is automatically reverted, structurally protecting consumers.

% =============================================================================
\section{Economic Model and Simulation}
\label{sec:economic}
% =============================================================================

\subsection{Asset Parameterization and Simulation}
The base unit of analysis is a 450~W monocrystalline panel costing R6,750 (2024). Table~\ref{tab:params} summarizes the parameters.

\begin{table}[h]
\caption{Solar Asset Economic Parameters (450~W Panel)}
\label{tab:params}
\begin{center}
\renewcommand{\arraystretch}{1.2}
\begin{tabular}{lrl}
\toprule
\textbf{Parameter} & \textbf{Value} & \textbf{Unit / Note} \\
\midrule
Panel capacity    & 450            & W \\
Capital cost      & R6,750         & ZAR (2024) \\
Avg.\ sun-hours   & 5.5            & h~day$^{-1}$ (annual avg.) \\
Annual generation & 903.4          & kWh~year$^{-1}$ \\
EaaS rate         & $\leq$R2.23    & ZAR/kWh (10\% below Eskom) \\
Gross annual rev. & R2,014         & ZAR (at R2.23/kWh) \\
Net yield (EaaS)  & 18.7\%         & adjusted per annum \\
\bottomrule
\end{tabular}
\end{center}
\end{table}

An end-to-end simulation of the smart contracts on Optimism Sepolia validated the flow. The system ingested simulated TPM telemetry, passed the data through the yield engine, and triggered the \texttt{YieldDistributor} contract. A gross yield test of R34,223.75 successfully triggered the 75/8/7/5/5 waterfall, allocating exactly R31,656.97 (the 75\% investor tranche + reserves) in USDC-wei precision. 

While the protocol manages the ZAR-denominated tariff off-chain and distributes USDC on-chain, the \textbf{investor} inherently carries the Foreign Exchange (FX) risk. Depreciation of the ZAR against the USD reduces the effective yield.

\subsection{Sensitivity Analysis}
Real-world conditions introduce variance. Table~\ref{tab:sensitivity} models the impact of irradiance, panel degradation (0.5\%--0.7\%) \cite{vartiainen2020lcoe}, inverter loss (3\%--5\%), and most critically, \emph{collection risk} (the probability of a household failing to pay).

\begin{table}[h]
\caption{Sensitivity Analysis: Investor Yield Under Varying Assumptions}
\label{tab:sensitivity}
\begin{center}
\renewcommand{\arraystretch}{1.2}
\begin{tabular}{p{2.2cm}rrrrr}
\toprule
\textbf{Scenario} & \textbf{Sun-h} & \textbf{Deg.} & \textbf{Inv.} & \textbf{Coll.} & \textbf{Yield} \\
\midrule
Conservative     & 4.0 & 0.7\% & 5\% & 80\% & 11.0\% \\
Realistic (c=80\%) & 5.5 & 0.5\% & 4\% & 80\% & 14.9\% \\
Realistic (c=90\%) & 5.5 & 0.5\% & 4\% & 90\% & 16.8\% \\
Realistic (c=100\%)& 5.5 & 0.5\% & 4\% & 100\% & 18.7\% \\
High irradiance  & 6.5 & 0.3\% & 3\% & 100\% & 23.1\% \\
\bottomrule
\end{tabular}
\end{center}
\end{table}

The baseline adjusted yield is 18.7\%. However, if collection rates drop to 80\% (common in emerging market EaaS models), the yield drops to 14.9\%. 

\subsection{Gas Cost and Scalability}
The \texttt{YieldDistributor} utilizes a push-based model on Optimism (L2), which iterates over all token holders. At 0.001 Gwei, distributing to 1,000 holders costs approximately \$0.078. While economically viable, the Optimism block gas limit imposes a hard ceiling of $\approx$1,150 holders per transaction. For larger deployments, the architecture dictates a migration to a Merkle-proof pull distribution pattern.

% =============================================================================
\section{Generalizability and Future Capabilities}
\label{sec:generalizability}
% =============================================================================

While instantiated for solar PV, the Cryptoeconomic Lifecycle Architecture is asset-agnostic. It applies equally to any DePIN asset generating cash flows based on physically verifiable work.

A primary future capability lies in \textbf{pooling excess battery storage} for high-demand third parties. As distributed solar installations increasingly pair with local battery storage, excess generation can be aggregated into a virtual power plant (VPP). Rather than solely servicing the local household, this pooled storage can be sold directly to data centers, whose spatio-temporal flexibility and massive energy demand make them ideal buyers in regional P2P energy markets \cite{jin2025data}. 

Under our framework, this simply extends the lifecycle: the battery management system utilizes the same Stage 3 TPM hardware to sign export telemetry. The Stage 4 settlement smart contract algorithmically routes data center revenue back to the fractional panel owners. By formalizing the trust boundaries, the architecture supports complex multi-party energy trading without requiring ad-hoc trust assumptions for each new asset class.

% =============================================================================
\section{Security and Trust Model Evaluation}
\label{sec:evaluation}
% =============================================================================

\subsection{The Two-Tier Trust Model}
To accurately assess the Veridian prototype, we apply a \emph{two-tier trust model} distinguishing between design-level guarantees and evaluation-level implementation. 

\textbf{Tier 1: Hardware Attestation.} At the design level, tampering requires physical extraction of a non-exportable TPM private key. However, at the evaluation level, the current prototype uses software-simulated keys. The security guarantee is therefore currently unverifiable until physical hardware is deployed.

\textbf{Tier 2: Oracle Delivery.} At the design level, yield distribution is triggered by a decentralized consensus network (e.g., Chainlink). At the evaluation level, it is currently triggered by a centralized FastAPI backend. 

\subsection{Comparison with Existing Protocols}
Table~\ref{tab:comparison} demonstrates how an explicit trust mapping differentiates this architecture from existing models.

\begin{table}[h]
\caption{Comparison with Related Energy and RWA Protocols}
\label{tab:comparison}
\renewcommand{\arraystretch}{1.2}
\begin{tabular}{p{1.6cm}p{1.7cm}p{1.6cm}p{2.1cm}}
\toprule
\textbf{Feature} & \textbf{Veridian} & \textbf{Power Ledger} & \textbf{Sun Exchange} \\
\midrule
Attestation & TPM 2.0 (Sim) & Grid Meter & Off-chain Trust \\
Ownership & ERC-3643 & Utility Token & Centralized DB \\
Settlement & USDC (L2) & POWR & BTC / ZAR \\
Protection & On-chain Cap & Market Rate & None \\
\bottomrule
\end{tabular}
\end{table}

Unlike Power Ledger, which assumes functioning grid interconnection and relies on unverified grid meters, Veridian anchors trust in local hardware cryptography. Unlike Sun Exchange, which maintains centralized custody of ownership ledgers and yield distributions, Veridian enforces compliance via ERC-3643 while removing custodial risk from the settlement stage.

% =============================================================================
\section{Conclusion}
\label{sec:conclusion}
% =============================================================================

This paper introduced a generalized Cryptoeconomic Lifecycle Architecture for decentralized energy assets, mapping the explicit trust boundaries across Capital Formation, Compliance, Telemetry, Settlement, and Governance. Through the Veridian solar EaaS prototype, we demonstrated that physical telemetry can be cryptographically linked to regulated on-chain settlement, producing a verifiable 18.7\% yield for retail investors. 

By providing an honest, rigorous map of where trust relies on cryptography versus law or plain faith, this architecture offers a blueprint for DePIN protocols to scale out of permissionless sandboxes and into complex, regulated infrastructure financing environments.

% =============================================================================
\section*{Acknowledgment}

The authors thank the faculty and teaching assistants of the UZH Deep Dive into Blockchain 2026 programme, and the UZH Blockchain Center for providing the research environment and review support that enabled this work.

% =============================================================================
\bibliographystyle{IEEEtran}
\bibliography{veridian_ieee}

\end{document}
